\documentclass[11pt,a4paper]{article}
\usepackage[utf8]{inputenc}
\usepackage[T1]{fontenc}
\usepackage{xeCJK}
\usepackage{amsmath, amssymb}
\usepackage{graphicx}
\usepackage{booktabs}
\usepackage{fullpage}
\usepackage{hyperref}
\usepackage{float}
\usepackage{caption}

\setCJKmainfont{Noto Serif CJK SC}

\title{\textbf{BVR-Sim —— 针对现有超视距空战强化学习环境的改进}}
\author{孙浩程 \\ 北京邮电大学}
\date{}

\begin{document}
\maketitle

\section{研究背景与目标}

Beyond-Visual-Range（BVR）空战是现代空战中的关键作战形态，涉及多智能体博弈、雷达感知约束以及导弹动力学控制等复杂问题。当前在该领域的强化学习（Reinforcement Learning, RL）研究中，存在以下两个主要困难：

\begin{itemize}
    \item \textbf{缺乏高保真且可复现的仿真环境：} 军用仿真平台通常闭源不可用，而开源环境往往简化了空气动力学、导弹制导或感知模型；
    \item \textbf{难以实现可对比的算法研究：} 不同研究中使用的观测、奖励、数据记录机制不统一，导致实验结果难以复现或比较。
\end{itemize}

为此，本课题旨在设计并实现一个开源、高保真、强化学习友好的远程空战仿真环境——BVR-Sim。  
该环境目标是兼顾物理真实性与强化学习实验可用性，为研究者提供统一的实验接口，用于研究：
\begin{itemize}
    \item 多智能体强化学习（MARL）；
    \item 自博弈（self-play）策略演化；
    \item 策略蒸馏与模仿学习（distillation / imitation learning）；
    \item 离线强化学习（offline RL）与数据驱动策略优化。
\end{itemize}

% \section{系统设计与主要改进}

% BVR-Sim 的设计遵循三个原则：物理建模相对真实、动作空间抽象且可通用、可视化与记录完备。  
% 与现有环境（如 BVR-Gym、B-ACE、WUKONG）相比，BVR-Sim 不仅在动力学上更高保真，还在系统交互、奖励结构与多智能体训练支持上进行了创新。

% \subsection{系统架构与模块化设计}

% 环境采用模块化的 Mission Runner 结构，包括：
% \begin{itemize}
%     \item \textbf{Physics：} 飞机与导弹的动力学求解；
%     \item \textbf{Sensing：} 雷达与数据链建模；
%     \item \textbf{Recorder：} 轨迹记录与离线数据集生成；
%     \item \textbf{Visualizer：} 奖励项与状态可视化；
% \end{itemize}

% 系统完全由 JSONC 配置文件驱动，可快速修改场景参数（阵位、目标数量、奖励权重等）而无需修改代码。

% \subsection{物理与导弹动力学建模改进}

% BVR-Sim 对导弹气动特性进行了高保真建模，采用基于 Mach 数的阻力模型（Mach-indexed drag）：
% \[
% D = \frac{1}{2}\rho(h)v^2 S_{\text{ref}} C_x(M,\alpha)
% \]
% 其中阻力系数 $C_x(M,\alpha)$ 含有跨音速波阻与超音速修正项，更符合真实导弹包线。  
% 此外，动力学模型同步包含推力、重力、气动损失等因素，并保证积分步长与 RL 时间步对齐，从而保证训练的物理一致性。

% \subsection{混合制导律（Hybrid Guidance）}

% 在制导模型上，引入了角度误差与视线角速率的混合控制律，使推力段和滑翔段具备不同控制行为：
% \[
% \dot{\beta}_{\text{cmd}} = \lambda(t)k_\beta[\phi-(\beta-\pi)] + (1-\lambda(t))\dot{\beta}_{\text{meas}}
% \]
% 其中 $\lambda(t)$ 为时间加权函数，在发动机推力期内逐步从角度控制过渡到速率控制，显著提升了制导稳定性与真实性。

% \subsection{Reward composition and visualization}
% Rewards are decomposed into dense shaping and sparse event terms:
% \begin{align}
% R_t &= \sum_{c\in\mathcal{C}} w_c r^c_t, \qquad \mathcal{C}=\{\text{engage},\text{altitude},\text{safe-alt},\text{launch},\text{hit},\text{win},\ldots\}, \label{eq:reward_sum}
% \end{align}
% with weights configured via JSONC. A per-component visualizer emits JSON and PNG plots (symlog scale) each episode, enabling diagnosis of mode collapse or overly sparse signals.

% \subsection{空战态势评估与观测向量扩展}

% 为更准确刻画空战过程，BVR-Sim 引入了态势评估指标（Situation Awareness Metrics），如能量比、角度优势、导弹威胁指数：
% \[
% S_{\text{threat}} = \frac{1}{1+\exp(-k(d_{\text{missile}}-d_0))}
% \]
% 这些指标可被直接纳入观测空间，或在评估中用于分析智能体决策质量。  
% 此外，观测噪声遵循可配置的高斯分布模型：
% \[
% \tilde{\mathbf{p}} = \mathbf{p} + \mathcal{N}(0,\sigma_p^2I_3)
% \]
% 用于模拟雷达测距误差与信息不完备性。

% \subsection{自博弈与训练系统支持}

% BVR-Sim 设计了轻量级的自博弈框架（Self-Play Framework）：
% \begin{itemize}
%     \item 支持基于 Elo 评分的策略对战选择；
%     \item 提供离线对战缓存与轨迹回放机制；
%     \item 可直接用于多智能体算法（如 MAPPO、HAPPO）训练；
%     \item 内置对称阵位生成器（Symmetric Initializer）以增强策略泛化。
% \end{itemize}

% \subsection{可视化与调试工具链}

% BVR-Sim 内置多项辅助工具：
% \begin{itemize}
%     \item \textbf{Reward Visualizer：} 自动生成每个奖励项的趋势图（JSON + PNG），用于诊断模式塌陷与稀疏奖励问题；
%     \item \textbf{ACMI 导出与 Trajectory Recorder：} 可在 Tacview 或自制工具中重放；
%     \item \textbf{Profiler 模块：} 记录各阶段耗时并计算统计量：
%     \[
%     \mu_k = \frac{1}{N_k}\sum_{i=1}^{N_k}\Delta t_{k,i}, \quad
%     \sigma_k = \sqrt{\frac{1}{N_k-1}\sum_{i=1}^{N_k}(\Delta t_{k,i}-\mu_k)^2}
%     \]
% \end{itemize}
\section{系统设计与主要改进}

BVR-Sim 的系统设计从三个核心目标出发：物理真实性、强化学习可用性以及研究可复现性。具体改进与特性如下：

\subsection{物理与制导建模}

\paragraph{飞机与导弹动力学}
环境提供可插拔的飞行器动力学求解器：
\begin{align}
\mathbf{x}_{t+1} &= \mathbf{x}_t + f_{\text{FDM}}\big(\mathbf{x}_t, \mathbf{u}_t(\Delta\psi_t,\Delta h_t,\Delta v_t)\big)\Delta t,
\end{align}
其中 $\mathbf{x}$ 包括位置、速度与姿态，$\mathbf{u}_t$ 是基于离散动作映射的自律率指令：
\begin{align}
\Delta \psi_t &= \gamma_\psi \frac{k_\psi-\bar{k}_\psi}{\bar{k}_\psi}, &
\Delta h_t &= \gamma_h \frac{k_h-\bar{k}_h}{\bar{k}_h}, &
\Delta v_t &= \gamma_v \frac{k_v-\bar{k}_v}{\bar{k}_v}.
\end{align}

导弹动力学考虑推力、Mach 阻力、重力及气动损失，并同步于飞机求解器：
\begin{align}
D &= \tfrac{1}{2}\rho(h)v^2 S_{\text{ref}} C_x(M,\alpha), \\
C_x(M,\alpha) &= C_{x0}(M) + C_{xB}(M) + K_1(M)|\alpha| + K_2(M)\alpha^2 + C_{x,\text{wave}}(M) + C_{x,\text{sup}}(M).
\end{align}
此设计允许导弹阻力随 Mach 数和攻角变化，真实再现跨音速、超音速性能。

\paragraph{混合制导模型}
导弹采用角误差与视线率（LOS-rate）混合控制：
\begin{align}
\dot{\beta}_{\text{cmd}} &= \lambda(t)\,k_\beta\big[\phi - (\beta - \pi)\big] + (1-\lambda(t))\,\dot{\beta}_{\text{meas}}, \\
n_y &= \frac{K(R)}{g}\,v\cos\theta\,\dot{\beta}_{\text{cmd}}, &
n_z = \frac{K(R)}{g}\,v\,\dot{\epsilon} + \cos\theta.
\end{align}
负载因子 $n_y,n_z$ 硬性设置了最大值，模拟导弹最大机动性能。融合了角度和角速度的引导率更接近现实中先进空空导弹的引导率，能飞出更真实的弹道。

\subsection{动作与观测接口}
动作空间采用Multi-Discrete编码映射到自律率指令，并增加导弹发射二值开关：
\begin{align}
\mathbf{a}_t = (k_\psi, k_h, k_v, k_\ell), \quad k_\ell\in\{0,1\}.
\end{align}

观测空间可配置（紧凑/扩展/信息不完全模式），包括自身、友方、敌方以及飞行中导弹的归一化相对运动状态：
\begin{align}
o_t &= \big[s^{\text{self}}_t,\ \{s^{\text{ally}}_{t,i}\}_{i=1}^{N_a},\ \{s^{\text{enemy}}_{t,j}\}_{j=1}^{N_e},\ \{s^{\text{miss}}_{t,m}\}_{m=1}^{M}\big], \\
\tilde{\mathbf{p}}_{t} &= \mathbf{p}_{t} + \boldsymbol{\epsilon}_p, & \tilde{\mathbf{v}}_{t} &= \mathbf{v}_{t} + \boldsymbol{\epsilon}_v, \\
\boldsymbol{\epsilon}_p &\sim \mathcal{N}(0,\sigma_p^2 I_3), \quad \boldsymbol{\epsilon}_v \sim \mathcal{N}(0,\sigma_v^2 I_3).
\end{align}

\subsection{训练工具链与可复现性}

\paragraph{奖励分解与可视化}  
奖励组合包括稠密引导与稀疏事件：
\begin{align}
R_t &= \sum_{c\in\mathcal{C}} w_c r^c_t, \quad
\mathcal{C} = \{\text{engage},\text{altitude},\text{safe-alt},\text{launch},\text{hit},\text{win},\ldots\}.
\end{align}
每个组件可生成 JSON 与 PNG 可视化，帮助分析训练信号稀疏性或reward数量级不匹配等问题。

\paragraph{蒸馏与教师强制}  
支持基于策略蒸馏或教师强制的训练：
\begin{align}
\mathcal{L}_{\text{dist},t} &= \lVert \hat{\mathbf{a}}^{\text{RL}}_t - \hat{\mathbf{a}}^{\star}_t\rVert_p + w_{\text{shoot}}|\sigma^{\text{RL}}_t - \sigma^{\star}_t|, \\
\tilde{\mathbf{a}}_t &= (1-\eta_t)\mathbf{a}^{\text{RL}}_t + \eta_t\mathbf{a}^{\star}_t, \quad \eta_t\in\{0,1\}.
\end{align}

\paragraph{轨迹记录与离线 RL}  
环境记录飞机与导弹状态轨迹：
\begin{align}
(o_t, \mathbf{a}_t, r_t, o_{t+1}, \text{aux}_t),
\end{align}
可用于离线强化学习与行为克隆，并支持批量存储与安全归档。

\subsection{系统模块化与可扩展性}
BVR-Sim 模块化设计：
\begin{itemize}
    \item 飞行器、导弹、传感器模型可替换；
    \item 动作接口与观测接口可自由配置；
    \item JSONC 配置驱动场景生成、奖励权重与日志路径；
    \item 自博弈、多团队训练框架内置。
\end{itemize}
这种设计兼顾高保真仿真与 RL 实验可用性，实现环境灵活性与研究可复现性的统一。

\section{总体效果与优势（总结与展望）}

BVR-Sim 的总体优势可总结为以下几点：
\begin{itemize}
    \item \textbf{高保真物理：} 同时考虑 Mach 阻力、重力、能量约束和混合制导；
    \item \textbf{强化学习可用性：} 多离散动作设计与可配置观测空间适配主流 RL 算法；
    \item \textbf{研究可复现性：} 内置奖励分解、轨迹记录、性能统计；
    \item \textbf{可扩展性：} 模块化结构支持替换不同的飞行器、导弹或传感器模型；
    \item \textbf{训练系统化：} 提供自博弈框架与标准化奖励接口，便于算法比较。
\end{itemize}

未来工作包括扩展更多飞行器与导弹模型、增加编队与协同任务支持、并发布标准化数据集与基准测试，以进一步推动空战强化学习的可复现研究。

\begin{thebibliography}{9}
\bibitem{scukins2024bvrgym}
E.~Scukins, M.~Klein, L.~Kroon, and P.~\"Ogren, ``BVR Gym: A Reinforcement Learning Environment for Beyond-Visual-Range Air Combat,'' \emph{arXiv preprint} arXiv:2403.17533, 2024.

\bibitem{kuroswiski2024bace}
A.~R. Kuroswiski, A.~S. Wu, and A.~Passaro, ``B-ACE: An Open Lightweight Beyond Visual Range Air Combat Simulation Environment for Multi-Agent Reinforcement Learning,'' 2024. DOI: 10.13140/RG.2.2.11999.57762.

\bibitem{piao2020wukong}
H.~Piao \emph{et~al.}, ``WUKONG: Beyond-Visual-Range Air Combat Tactics Auto-Generation by Reinforcement Learning,'' \emph{IEEE}, 2020.
\end{thebibliography}

\end{document}
